\subsection{Selection of Bad Epochs}

The first stage of the proposed pipeline performs epoch-level screening using an
autoreject-inspired threshold selection procedure. Its purpose is not to estimate
the onset or offset of a blink, but to identify epochs whose amplitudes are
sufficiently abnormal to warrant subsequent sample-level localization. Let
\(X \in \mathbb{R}^{N \times C \times T}\) denote the epoched multichannel
signal after preprocessing, where \(N\) is the number of valid epochs, \(C\) is
the number of channels, and \(T\) is the number of samples per epoch. The sample
at epoch \(i \in \mathcal{E}=\{1,\ldots,N\}\), channel
\(c \in \mathcal{C}=\{1,\ldots,C\}\), and within-epoch time index
\(t \in \{1,\ldots,T\}\) is denoted by \(x_{i,c,t}\).

For each epoch and channel, an amplitude-range statistic is computed using the
peak-to-peak value
\[
a_{i,c}
= \max_{1 \leq t \leq T} x_{i,c,t}
- \min_{1 \leq t \leq T} x_{i,c,t}.
\]
Large values of \(a_{i,c}\) indicate that the corresponding epoch contains a
high-amplitude excursion in channel \(c\), a pattern consistent with ocular or
other transient artifacts. Rather than imposing a fixed rejection level, the
pipeline estimates a channel-specific threshold \(\tau_c^\star\) from the data.
Candidate thresholds are taken from the empirical range of peak-to-peak
amplitudes for that channel, denoted
\(\mathcal{Q}_c=\{a_{i,c}: i \in \mathcal{E}\}\), and are evaluated by a
cross-validation criterion.

Specifically, let \(\{(\mathcal{T}_k,\mathcal{V}_k)\}_{k=1}^{K}\) be
cross-validation splits of the epoch indices into training sets
\(\mathcal{T}_k\) and validation sets \(\mathcal{V}_k\). For a candidate
threshold \(\tau \in \mathcal{Q}_c\), the training epochs retained for channel
\(c\) in fold \(k\) are
\[
\mathcal{G}_{k,c}(\tau)
= \{ i \in \mathcal{T}_k : a_{i,c} \leq \tau \}.
\]
These retained epochs define a threshold-dependent training template
\[
\bar{x}^{(\tau)}_{k,c}(t)
= \frac{1}{|\mathcal{G}_{k,c}(\tau)|}
\sum_{i \in \mathcal{G}_{k,c}(\tau)} x_{i,c,t},
\]
where candidates yielding an empty retained set are treated as inadmissible. The
corresponding validation reference is computed as the pointwise median across
the validation epochs,
\[
\tilde{x}_{k,c}(t)
= \operatorname{median}_{i \in \mathcal{V}_k} x_{i,c,t}.
\]
The use of a median validation reference makes the criterion less sensitive to
large validation-set artifacts. The validation loss for channel \(c\) is then
defined as
\[
\mathcal{J}_{c}(\tau)
= \frac{1}{K}\sum_{k=1}^{K}
\left[
\frac{1}{T}\sum_{t=1}^{T}
\left\{\bar{x}^{(\tau)}_{k,c}(t)-\tilde{x}_{k,c}(t)\right\}^{2}
\right]^{1/2}.
\]
The selected channel threshold is the candidate that minimizes the average
validation discrepancy,
\[
\tau_c^\star
= \arg\min_{\tau \in \mathcal{Q}_c} \mathcal{J}_{c}(\tau).
\]
In practice, this objective may be searched by sequential Bayesian optimization
or by a random search over candidate peak-to-peak thresholds; both choices
retain the same validation objective and rejection rule.

After the thresholds \(\{\tau_c^\star\}_{c=1}^{C}\) have been estimated, each
epoch receives a channel-level artifact indicator
\[
z_{i,c}
= \mathbf{1}\{a_{i,c} > \tau_c^\star\},
\]
where \(\mathbf{1}\{\cdot\}\) is the indicator function. The present pipeline
uses the union of the channel-level decisions to form the epoch-level screening
mask,
\[
b_i
= \mathbf{1}\left\{\sum_{c=1}^{C} z_{i,c} \geq 1\right\}
= \mathbf{1}\left\{\exists c \in \mathcal{C}: a_{i,c} > \tau_c^\star\right\}.
\]
The set of suspicious epochs is therefore
\[
\mathcal{B}=\{i \in \mathcal{E}: b_i=1\}.
\]
These epochs are interpreted as likely blink-contaminated or artifact-heavy
segments. They provide the restricted data subset from which the next stage
estimates sample-level blink-region thresholds.
